% Part: first-order-logic
% Chapter: models-theories
% Section: expressing-relations

\documentclass[../../../include/open-logic-section]{subfiles}

\begin{document}

\olfileid{fol}{mat}{exr}

\olsection{Relaties uitdrukken in \article{structure}
  \printtoken{S}{structure}}

\begin{explain}
Een belangrijke toepassing van !!{formula}s is eigenschappen en
relaties in !!a{structure}~$\Struct M$ uit te drukken met de
primitieve begrippen van de taal~$\Lang L$ van~$\Struct M$. Daarmee
bedoelen we het volgende. Het !!{domain} van $\Struct M$ is een
verzameling objecten. De !!{constant}s, !!{function}s en
!!{predicate}s worden in~$\Struct M$ geïnterpreteerd door bepaalde
objecten in~$\Domain M$, functies op~$\Domain M$ en relaties
op~$\Domain M$. Als bijvoorbeeld $\Obj A^2_0$ in $\Lang L$ voorkomt,
kent $\Struct M$ daaraan een relatie~$R = \Assign{{\Obj
  A^2_0}}{M}$ toe. De formule $\Atom{\Obj A^2_0}{\Obj v_1, \Obj v_2}$
\emph{drukt dan precies die relatie uit} in de volgende zin:
als een bedeling~$s$ de variabele $\Obj v_1$ op $a \in \Domain{M}$ en
$\Obj v_2$ op $b \in \Domain M$ afbeeldt, dan
\[
Rab \text{\quad iff\quad} \Sat{M}{\Atom{\Obj A^2_0}{\Obj v_1, \Obj v_2}}[s].
\]
Merk op dat we hier bedelingen moeten gebruiken: we kunnen niet zonder
meer zeggen ``$Rab$ dan en slechts dan als
$\Sat{M}{\Atom{\Obj A^2_0}{a, b}}$'', want $a$ en $b$ zijn geen
symbolen van onze taal maar !!{element}s van~$\Domain{M}$.

We beschikken niet alleen over atomaire !!{formula}s, maar kunnen ze
met logische connectieven en kwantoren combineren. Daardoor kunnen
complexere !!{formula}s andere relaties definiëren die niet
rechtstreeks in~$\Struct M$ zijn ingebouwd. We onderzoeken hoe dat
gebeurt en in het bijzonder welke relaties in !!a{structure}
definieerbaar zijn.
\end{explain}

\begin{defn}
Zij $!A(\Obj v_1,\dots, \Obj v_n)$ !!a{formula} van $\Lang L$ waarin
alleen $\Obj v_1$,\dots, $\Obj v_n$ vrij voorkomen, en zij $\Struct M$
!!a{structure} voor~$\Lang L$. De formule
$!A(\Obj v_1,\dots, \Obj v_n)$ drukt de relatie~$R \subseteq
\Domain M^n$ \emph{uit} dan en slechts dan als
\[
Ra_1\dots a_n \text{\quad iff\quad} \Sat{M}{\Atom{!A}{\Obj
    v_1,\dots, \Obj v_n}}[s]
\]
voor iedere bedeling~$s$ waarvoor $s(\Obj v_i) = a_i$ ($i = 1,
\dots, n$).
\end{defn}

\begin{ex}
In het standaardmodel van de rekenkunde~$\Struct N$ drukt de
!!{formula} $\Obj v_1 < \Obj v_2 \lor \eq[\Obj v_1][\Obj v_2]$ de
relatie~$\le$ op~$\Nat$ uit. De !!{formula}
$\eq[\Obj v_2][\Obj v_1']$ drukt de opvolgerrelatie uit, dat wil zeggen
de relatie $R \subseteq \Nat^2$ waarvoor $Rnm$ geldt als $m$ de
opvolger van~$n$ is. De formule $\eq[\Obj v_1][\Obj v_2']$ drukt de
voorgangerrelatie uit. De !!{formula}s
$\lexists[\Obj v_3][(\eq/[\Obj v_3][\Obj 0] \land
  \eq[\Obj v_2][(\Obj v_1 + \Obj v_3)])]$ en $\lexists[\Obj
  v_3][\eq[(\Obj v_1 + {\Obj v_3}')][v_2]]$ drukken beide de relatie
$\Obj <$ uit. Het predicaatsymbool~$<$ is dus eigenlijk overbodig in
de taal van de rekenkunde: het kan worden gedefinieerd.
\end{ex}

\begin{explain}
Deze gedachte is niet alleen van belang voor afzonderlijke
!!{structure}s, maar telkens wanneer we met een taal één of meer
beoogde modellen beschrijven, dus wanneer we theorieën beschouwen.
Zulke theorieën bevatten vaak slechts enkele !!{predicate}s als
basissymbolen, terwijl in het beschreven domein veel andere relaties
een belangrijke rol spelen. Als die andere relaties stelselmatig
kunnen worden uitgedrukt met de relaties die de basis-!!{predicate}s
van de taal interpreteren, zeggen we dat ze in de taal kunnen worden
\emph{gedefinieerd}.
\end{explain}

\begin{prob}
Vind !!{formula}s in $\Lang L_A$ die de volgende relaties definiëren:
\begin{enumerate}
\item $n$ ligt tussen $i$ en $j$;
\item $n$ deelt $m$ (dus $m$ is een veelvoud van $n$);
\item $n$ is een priemgetal (dus geen enkel getal behalve $1$ en $n$
  deelt~$n$).
\end{enumerate}
\end{prob}

\begin{prob}
Veronderstel dat de formule $!A(\Obj v_1, \Obj v_2)$ de relatie $R
\subseteq \Domain M^2$ in !!a{structure}~$\Struct M$ uitdrukt. Vind
formules die de volgende relaties uitdrukken:
\begin{enumerate}
\item de inverse $R^{-1}$ van $R$;
\item de samenstelling (het relatieve product) $R \mid R$;
\end{enumerate}
Kunt u ook $R^+$, de transitieve afsluiting van~$R$, uitdrukken?
\end{prob}

\begin{prob}
Zij $\Lang{L}$ de taal die alleen het tweeplaatsige predicaatsymbool
$<$ bevat (dus geen andere !!{constant}s, !!{function}s of
!!{predicate}s---behalve uiteraard~$\eq$). Zij $\Struct{N}$ de
structuur waarvoor $\Domain{N} = \Nat$ en
$\Assign{<}{N} = \Setabs{\tuple{n,m}}{n < m}$. Bewijs het volgende:
\begin{enumerate}
\item $\{ 0 \}$ is definieerbaar in $\Struct{N}$;
\item $\{ 1 \}$ is definieerbaar in $\Struct{N}$;
\item $\{ 2 \}$ is definieerbaar in $\Struct{N}$;
\item voor iedere $n \in \Nat$ is de verzameling $\{ n \}$
  definieerbaar in $\Struct{N}$;
\item iedere eindige deelverzameling van $\Domain{N}$ is definieerbaar
  in $\Struct{N}$;
\item iedere co-eindige deelverzameling van $\Domain{N}$ is
  definieerbaar in $\Struct{N}$ (waarbij $X \subseteq \Nat$
  co-eindig is dan en slechts dan als $\Nat \setminus X$ eindig is).
\end{enumerate}
\end{prob}


\end{document}
